\chapter{Метод прийняття рішень контекстної безпеки з урахуванням політик}
\label{ch:policy}

% ============================================================
%  РОЗДІЛ 4 -- другий центральний метод середовища виконання політик AURA.
%  Перетворює сигнал ризику (розділ 3) у захисну дію.
% ============================================================

\section{Перетворення сигналу ризику в захисну дію}
\label{sec:signal-to-action}
% -- Від мітки/оцінки ризику до рішення; принцип урахування політик.

\section{Простір рішень системи обміну повідомленнями}
\label{sec:decision-surface}
% -- Дії: дозвіл, попередження, ескалація, карантин, блокування,
%    передавання на контрольований перегляд (зв'язок з розділом 7).

\section{Опрацювання невизначеності}
\label{sec:uncertainty}
% -- межі безпечного закриття й контрольованого пропуску; пороги; стан утримання від рішення.

\section{Профіль захисту учасника та доменний режим}
\label{sec:protection-profile}
% -- Налаштування рішень за профілем (зокрема дитячий) і доменом.

\section{Приватний аудиторський запис}
\label{sec:audit-record}
% -- Що фіксується без розкриття приватного змісту; мінімізація.

\section{Критерії якості з контрольованим допуском до випуску}
\label{sec:release-gates}
% -- Статуси PASS / FAIL / INSUFFICIENT_SUPPORT / BLOCKED (деталі -- розділ 8).
Прийняття рішень доповнюється критеріями якості з контрольованим допуском до випуску зі статусами \texttt{PASS}, \texttt{FAIL}, \texttt{INSUFFICIENT\_SUPPORT} та \texttt{BLOCKED}, що формалізують умови допуску захисної логіки до використання (методологію оцінювання розгорнуто в розділі~8). \emph{[Розкрити критерії.]}

\rozdilconcl{4}
\emph{[Висновки до розділу 4.]}
